\documentclass{article}[12pt, a4paper]
\title{Electrostatics}
\usepackage{amsmath}
\usepackage{esint} % Required for \oiint and \oiiint
\usepackage{hyperref} %document outline
\usepackage{fancyhdr}
% \usepackage{lipsum} % For dummy text
\usepackage{xcolor}
% Page style setup
\pagestyle{fancy}
\fancyhf{} % Clear headers and footers
\renewcommand{\headrulewidth}{0pt} % <-- THIS REMOVES THE ANNOYING TOP LINE
\fancyfoot[R]{\textbf{\color{black}em-set-03}}
\newcommand{\gfrac}[2]{\frac{\displaystyle #1}{\displaystyle #2}}


\begin{document}
\centering \textbf{(Potential Problems)}
\begin{enumerate}


% Q1
\item A rectangular pipe, running parallel to the $z$-axis (from $-\infty$ to $+\infty$),
has three grounded metal sides, at $y = 0, y = a$, and $x = 0$. The fourth side, at
$x = b$, is maintained at a specified potential $V_0(y)$.
\begin{enumerate}
	\item Develop general formula for the potential inside the pipe.
    \item Find the potential explicitely for the, $V(y) = V_o$ (Constant).
\end{enumerate}

% Q2
\item A spherical shell with given potential $V_o(R, \theta) = kcos3\theta$ (at radius $r = R$, where $k$ is a constant).
\begin{enumerate}
	\item Find the potential inside the shell.
	\item Find the potential outside the shell.
	\item Calculate the total induced surface charge density $\sigma(\theta)$ in the shell.
\end{enumerate}

% Q3
\item Suppose the potential $V(\theta)$ at the surface of a sphere is specified,
and there is no charge inside or outside the sphere. Show that the charge density on
the sphere is given by,
$$
\sigma(\theta) = \gfrac{\epsilon_o}{2R}\sum_{l = 0}^{\infty}(2l + 1)^2C_lP_l(\cos\theta)
$$ where
$$
C_l = \int^{\pi}_{0} V_o(\theta) P_l(\cos\theta) \sin(\theta) d\theta$$


\end{enumerate}
\end{document}
